\section{Limitations and validity boundaries}
\label{sec:limitations}

\textbf{Scale and replication.} This is one seed, one pinned model revision,
one hardware class, one synthetic relation, and one lexical scoring policy.
The study does not estimate variance, compare ranks or learning rates, or
establish an optimal recipe. Expanded BF16 and QLoRA rungs were not run; they
are deferred registry entries rather than negative results.

\textbf{Evaluation scope.} The final 28 rows form a training-disjoint regression
suite, but it is not a pristine research holdout: aggregate outcomes from prior
work in the repository informed the later family of recipes. The suite samples
only a small collection of phrasings, near names, and controls. Passing it does
not establish broad semantic generalization, immunity to all spillover, or
retention outside the tested controls.

\textbf{Interpretation.} Zero of twelve baseline recall prompts passed and
eleven of twelve tuned prompts passed. This supports the registered
\Code{candidate-knowledge-acquisition} label under this protocol. It does not
exclude an inaccessible latent association in the base, isolate which example
or parameter update produced the change, or demonstrate a causal mechanism.
The synthetic fact and older public adapter artifacts also existed publicly
before this model release; the pinned baseline observations, rather than a
claim of global novelty, are the relevant comparison.

\textbf{Verification.} The anonymous public check was a single-prompt smoke
test. Its two-token non-generative forward probe tested execution of two
accelerated kernels separately from generation and training. Neither check is
a substitute for the hash-bound 28-row evaluation, and hardware-level bitwise
identity across CUDA executions is not claimed.

\textbf{Evidence capture.} The retrieved 15-file allowlist was bound at
retrieval time. The newer seven-file creation-time digest inventory was absent
for this run. Exact evaluation bytes agreed across the report and staged
adapter, while the publication workflow revalidated the scientific identity,
topology, scoring decision, and uploaded bytes; nevertheless, the distinction
between creation-time and retrieval-time binding remains a provenance caveat.
\claimsource{metadata}\claimsource{publication}

\textbf{Authorship.} LLMs substantially assisted the work, as disclosed on the
first page. Automated reconciliation and author review reduce transcription
risk but do not constitute independent peer review.
